\documentclass[UTF8,11pt]{ctexart}
\usepackage[a4paper,margin=24mm]{geometry}
\usepackage{amsmath,amssymb,amsthm,mathtools}
\usepackage[all]{xy}
\usepackage{xurl,hyperref,enumitem}
\hypersetup{hidelinks,breaklinks=true}
\setlength{\emergencystretch}{3em}
\newtheorem*{theorem}{Theorem}
\newtheorem*{lemma}{Lemma}
\newtheorem*{proposition}{Proposition}
\newtheorem*{corollary}{Corollary}
\theoremstyle{definition}
\newtheorem*{definition}{Definition}
\newtheorem*{remark}{Remark}
\DeclareMathOperator{\Hom}{Hom}
\DeclareMathOperator{\Ext}{Ext}
\DeclareMathOperator{\Tor}{Tor}
\DeclareMathOperator{\Spec}{Spec}
\DeclareMathOperator{\supp}{supp}
\DeclareMathOperator{\im}{im}
\DeclareMathOperator{\id}{id}
\DeclareMathOperator{\Tr}{Tr}
\DeclareMathOperator{\rank}{rank}
\newcommand{\R}{\mathbb R}
\newcommand{\C}{\mathbb C}
\newcommand{\Z}{\mathbb Z}
\newcommand{\N}{\mathbb N}
\newcommand{\Q}{\mathbb Q}
\providecommand{\fint}{\mathop{\rlap{\hspace{0.18em}$-$}\!\int}}
\begin{document}
\section*{024\quad 奇异卷积的Calderon--Zygmund有界性}
\subsection*{来源与收录范围}
Lawrence D. Guth（授课），Jane Wang（原文件署名转录），MIT 18.156，\emph{Differential Analysis II: Partial Differential Equations and Fourier Analysis}。采用 MIT OpenCourseWare Spring 2016 页面提供、文件内日期为 2015-04-13、04-15、04-17 的 Lectures 26--28。主命题为 Lecture 26 Theorem 1；其四部分证明分布于 Lecture 26 pp.~1--3、Lecture 27 \S2 pp.~2--3、Lecture 28 p.~1。分解引理只作为本题支撑，不另计题目。获取日期：2026-09-28。
\par\url{https://ocw.mit.edu/courses/18-156-differential-analysis-ii-partial-differential-equations-and-fourier-analysis-spring-2016/9298e2e46787978653f8d4c9689cf3a5_MIT18_156S16_lec26.pdf}
\par\url{https://ocw.mit.edu/courses/18-156-differential-analysis-ii-partial-differential-equations-and-fourier-analysis-spring-2016/c956b03645a1d98c3e8ceb7b86fe996d_MIT18_156S16_lec27.pdf}
\par\url{https://ocw.mit.edu/courses/18-156-differential-analysis-ii-partial-differential-equations-and-fourier-analysis-spring-2016/770671b5e8c06ea0875ca5e06b354d13_MIT18_156S16_lec28.pdf}
\paragraph{记号说明（整理者，非原文引句）} $S_r$ 为半径 $r$ 的球面，$V_h(\lambda)=|\{x:|h(x)|>\lambda\}|$，$\fint_Q h=|Q|^{-1}\int_Qh$。这里用 $\fint$ 排原文的平均积分，$\lesssim$ 保留原文含维数等固定参数的隐常数约定。奇异卷积采用课程上下文中的主值意义；本题不增加任意分布核的结论。
\subsection*{作者命题与人类证明}
\begin{theorem}[Lecture 26, Theorem 1]
If $Tf=f*K$ on $\R^d$, $|K(x)|\lesssim |x|^{-d}$, $|\partial K(x)|\lesssim |x|^{-d-1}$, $\int_{S_r}K(x)=0$ for all $r$, then
\[\|Tf\|_p\lesssim\|f\|_p,\qquad 1<p<\infty.\]
\end{theorem}
\begin{proof}[Part I: $L^2$ bound (Lecture 26)]
By Plancherel, we have that
\[\|f*K\|_2=\|\widehat f\cdot\widehat K\|_2\leq\|\widehat K\|_\infty\|\widehat f\|_2=\|\widehat K\|_\infty\|f\|_2.\]
So it suffices to bound $\|\widehat K\|_\infty$. We have that
\[|\widehat K(\omega)|=\left|\int K(x)e^{-i\omega x}\,dx\right|\leq\sum_{j\in\Z}\left|\int_{2^{j-1}\leq|x|\leq2^j}K(x)e^{-i\omega x}\,dx\right|=:\sum_{j\in\Z}I_j.\]
We'll also let $A_j:=\{x:2^{j-1}\leq|x|\leq2^j\}$. For small $j$, those such that $|\omega\cdot2^j|\leq1$, we have from $\int_{S_r}K(x)=0$ that
\begin{align*}
|I_j|&=\left|\int_{A_j}K(x)(e^{-i\omega x}-1)\,dx\right|\\
&\leq|\omega\cdot2^j|\int_{A_j}|K(x)|\sim|\omega\cdot2^j|(2^j)^d(2^j)^{-d}\sim|\omega\cdot2^j|.
\end{align*}
But now, $\sum_{|\omega w^j|\leq1}I_j$ is the sum of exponentially decreasing terms, and is therefore $\lesssim1$. We also have to worry about what happens for large $j$. For $j$ such that $|\omega\cdot2^j|>1$, we can choose $\ell\in\{1,2,\ldots,d\}$ such that $|\omega_\ell|\gtrsim|\omega|$ and integrate by parts to get that
\begin{align*}
|I_j|&=\left|\int_{A_j}K(x)e^{-i\omega x}\,dx\right|\\
&=\left|\int_{A_j}\partial_\ell K\cdot\frac1{i\omega_\ell}e^{-i\omega x}\,dx+\int_{\partial A_j}Ke^{-i\omega x}\cdot\frac1{i\omega_\ell}\,dx\right|\\
&\leq\int_{A_j}|\partial K|\cdot\frac1{|\omega|}\,dx+\int_{\partial A_j}|K|\cdot\frac1{|\omega|}\\
&\lesssim|A_j|(2^j)^{-d-1}\frac1{|\omega|}+|\partial A_j|(2^j)^{-d}\frac1\omega
\sim\frac1{|2^j\omega|}.
\end{align*}
Again, $\sum_{|\omega2^j|>1}I_j$ is bounded by an exponentially decaying series, and so this sum and therefore the whole sum $\lesssim1$. This gives us the $L^2$ bound that we wanted. We note here that sometimes in the statement of Calderon Zygmund, the $L^2$ bound $\|Tf\|_2\lesssim\|f\|_2$ is taken to be a condition instead of $\int_{S_r}K(x)=0$.
\end{proof}
\paragraph{Part II: Weak $L^1$ bound (Lecture 26)}
We want to prove the statement
\[V_{Tf}(\lambda)\lesssim\|f\|_1\lambda^{-1}.\tag{1}\]
We will do this by breaking up the function $f$ into a small part and a ``balanced part''. Let us first show that a weak $L^1$ bound holds for ``small'' and ``balanced'' functions. We'll start with small functions.
\begin{lemma}[Lecture 26, Lemma 2]
If $\|f\|_\infty\leq10\lambda$, then (1) holds.
\end{lemma}
\begin{proof}
This follows from the $L^2$ estimate.
\[V_{Tf}(\lambda)\leq\|Tf\|_2^2\lambda^{-2}\lesssim\|f\|_2^2\lambda^{-2}\lesssim\|f\|_1\lambda^{-1}.\]
\end{proof}
\begin{lemma}[Lecture 26, Lemma 3]
If $b(x)$ is ``balanced for $\lambda$'', $\supp b\subset\text{cube }Q$, $\fint_Q|b|=\lambda$, $\int_Qb=0$, then $|Tb(x)|\leq\lambda\cdot\mu^{-d-1}$. Here, $\mu s$ is the distance from $x$ to $Q$ and $\mu\geq2$.
\end{lemma}
\begin{proof}
Note that
\[|Tb(x)|\leq\int_Q|b|\cdot|K(x-y|\,dy\sim(\mu\cdot s)^{-d}\int_Q|b|\sim\lambda\cdot\mu^{-d},\]
but we can do better than that. If $y_0$ is the center of $Q$, then have that
\[|Tb(x)|=\left|\int_Qb(y)K(x-y)\,dy\right|=\left|\int_Qb(y)(K(x-y)-K(x-y_0))\,dy\right|.\]
Now, since $K(x-y)-K(x-y_0)\lesssim s\cdot\max_{y\in Q}|\partial K(x-y)|\lesssim s\cdot(\mu s)^{-d-1}$, we have that
\[|Tb(x)|\lesssim s\cdot(\mu s)^{-d-1}\int|b(y)|\sim\mu^{-d-1}\cdot\lambda.\]
\end{proof}
\begin{lemma}[Lecture 26, Lemma 4]
If $b=\sum b_j$, $b_j$ balanced functions for $\lambda$, and each function $b_j$ is supported on $Q_j$ disjoint sets, then $V_{Tb}(\lambda)\lesssim\|b\|_1\lambda^{-1}$.
\end{lemma}
\begin{proof}
We have that $\|b\|_1\sim\lambda\sum_j|Q_j|$. Let $U:=\bigcup_j2Q_j$. Then, $|U|\lesssim\|b\|_1\lambda^{-1}$. So it suffices to check that $\|Tb\|_{L^1(\R^d\setminus U)}\lesssim\|b\|_1$, and for this it suffices to check that $\|Tb_j\|_{L^1(\R^d\setminus2Q_j)}\lesssim\|b_j\|_1$, since then we would have that
\begin{align*}
\|Tb\|_{L^1(\R^d\setminus U)}&\leq\sum_j\|Tb_j\|_{L^1(\R^d\setminus U)}\\
&\leq\sum_j\|Tb_j\|_{L^1(\R^d\setminus2Q_j)}\lesssim\sum_j\|b_j\|_1=\|f\|_1,
\end{align*}
since the $b_j$ have disjoint supports. But that $\|Tb_j\|_{L^1(\R^d\setminus2Q_j)}\lesssim\|b_j\|_1$ follows from integrating the last lemma.
\end{proof}
\paragraph{Decomposition and completion of Part II (Lecture 27, \S2)}
\begin{lemma}[Lecture 27, Lemma 1]
For $f\in C_c^0$, $\lambda>0$, we can decompose $f=b+s$, the sum of a balanced part and a small part, such that $\|b\|_1+\|s\|_1\lesssim\|f\|_1$ and $\|s\|_\infty\leq\lambda$, $b=\sum b_j$ where $b_j$ are balanced for $\lambda$ supported on disjoint $Q_j$ and
\[\fint_{Q_j}b_j\lesssim\fint_{Q_j}f\lesssim\lambda.\]
\end{lemma}
\begin{proof}
We're going to use a Calderon-Zygmund iterated stopping time algorithm to construct $Q_j$ and $b_j$. Start with a cubical grid in $\R^d$ of side length $s$ large and $\fint_Q|f|<\lambda$ in each cube.

[Call this point in the algorithm (A).] Now, consider each $Q$.
\begin{enumerate}[label=(\roman*)]
\item If $\fint_Q|f|<\lambda$, subdivide $Q$ into $2^d$ equally sized cubes and repeat this step (A) with each of the subdivided cubes.
\item If $\fint_Q|f|\geq\lambda$, add $Q$ to the list of balanced cubes, call it $Q_j$, and let
\[b_j=f\cdot\chi_{Q_j}-\fint_{Q_j}f.\]
Do not go back to (A) with this cube.
\end{enumerate}
The output of the algorithm is $\{Q_j\}$ and a function $b_j$ for each $Q_j$. Then, let $b=\sum b_j$, $s=f-b$.

We can make some observations now. First,
\[\lambda\leq\fint_{Q_j}|f|<2^d\lambda.\]
We also have some bound for $s$. If $x\notin\bigcup Q_j$, then $|s(x)|=|f(x)|\leq\lambda$. If $x\in Q_j$, then
\[|s(x)|=|f(x)-b_j(x)|=\left|\fint_{Q_j}f\right|\leq\fint_{Q_j}|f|\leq2^d\lambda,\]
and so we have that
\[\int_{\R^d\setminus\bigcup Q_j}|s|=\int_{\R^d\setminus\bigcup Q_j}|f|\leq\|f\|_{L^1}.\]
From this, we get that
\[\int_{\bigcup Q_j}|s|\leq\int_{\bigcup Q_j}|f|\leq\|f\|_{L^1},\]
so $\|s\|_1\leq\|f\|_1$. We also have bounds for the $b_j$:
\[\fint_{Q_j}|b_j|=\fint_{Q_j}\left|f-\fint_{Q_j}f\right|\leq2\int_{Q_j}|f|\]
and
\[\int_{Q_j}b_j=\int_{Q_j}f-\int_{Q_j}f=0.\]
\end{proof}
This lemma then helps us conclude part II of the proof of Calderon-Zygmund, since $V_{Tf}(2\lambda)\leq V_{Ts}(\lambda)+V_{Tb}(\lambda)$. By the $L^2$ bound $V_{Tf}(\lambda)\lesssim\|f\|_1\lambda^{-1}$. We also have that
\begin{align*}
V_{Tb}(\lambda)&\leq\left|\bigcup_j2Q_j\right|+\lambda^{-1}\int_{\R^d\setminus\bigcup2Q_j}|Tb|\\
&\lesssim\left|\bigcup_jQ_j\right|+\sum_j\lambda^{-1}\int_{\R^d\setminus2Q_j}|Tb_j|\\
&\lesssim\|f\|_1\lambda^{-1}+\lambda^{-1}\sum_j\|b_j\|_1
\lesssim\lambda^{-1}(\|s\|_1+\|b\|_1)\lesssim\lambda^{-1}\|f\|_1.
\end{align*}
\paragraph{Part III: Interpolation (Lecture 27)}
Since we have a weak $L^1$ bound and a strong $L^2$ bound, we can use Marcinkiewicz interpolation to get the bound $\|Tf\|_p\lesssim\|f\|_p$ for $1<p\leq2$.
\begin{proof}[Part IV: Duality (Lecture 28)]
Recall that we have already proven CZ for $1<p\leq2$. Now, let $p$ be such that $2<p<\infty$ and let $p'$ be the dual exponent (so $1<p<2$). Then, we have that
\begin{align*}
\|Tf\|_p&=\sup_{\|g\|_{L^{p'}}\leq1}\int Tf\cdot g\\
&=\sup\int\!\int f(y)K(x-y)\,dy\,g(x)\,dx\\
&=\sup\int f(y)\int K(x-y)g(x)\,dx\,dy\\
&=\sup\int f\cdot(\overline K*g)
\leq\|f\|_p\cdot\sup\|\overline T g\|_{L^{p'}}\lesssim\|f\|_p,
\end{align*}
by CZ for $p'$. We note here that $\overline K(x)$ is defined as $K(-x)$ and $\overline T g=g*\overline K$. This completes the proof the the Calderon-Zygmund theorem.
\end{proof}
\subsection*{整理说明（非作者正文）}
本题计一个奇异积分 $L^p$ 有界性问题，四部分核心证明及停止时间分解均已收录；不另计分解引理，不把后一讲只证明单侧的 Littlewood--Paley 结论并入。标准先修为 Plancherel、积分分部、Chebyshev 不等式、Marcinkiewicz 插值、$L^p$ 对偶和稠密延拓；本题仍承担乘子分尺度估计、坏立方体选择、消去带来的远场增益和两端范围的组织，难度判断据此而非课程名称。

Lecture 26 的三个正文页、Lecture 27 的第2--3页及 Lecture 28 第1页均实际查看页面图像。保留原稿的若干简写/排印：分部积分边界项未展开法向量；第一部分有 $w^j$ 字样；Lemma 4 末式写 $\|f\|_1$；分解中的 $b_j$ 公式只在其支撑立方体上写出，正文已说支撑于 $Q_j$；平均积分不等式右端原页少一平均横线；小部估计一行写 $V_{Tf}$；对偶段括号写 $1<p<2$，而同段其余位置使用 $p'$。这里没有静默把这些改成新的人类证明，也不据此作数学认证。核反射的横线并非复共轭。空间立方体图只是相应文字定义的例示，未复制；没有删除额外证明步骤。原 PDF 未能保存为本地文件，\texttt{sources/} 保存的是明确标明归属的转录摘录。
\subsection*{可修改的简短标注（非作者正文）}
$c$：奇异积分算子的各指数范数有界性

$a$：Fourier 变换与 Plancherel；二进环带；球面消去；分部积分；停止时间立方体分解；零均值远场估计；弱型分布函数；Marcinkiewicz 插值；对偶核

\texttt{cross\_theory\_status = yes}。奇异卷积先经 Fourier 表示取得 L2 界，再把任意输入按空间尺度拆成小部和零均值部，利用消去控制弱 L1，最后经插值和对偶得到全部指数范围。位置：Lectures 26--28 的 Parts I--IV。

全部标注均为非穷尽、可修改初稿；未标注不表示未使用。
\subsection*{许可与修改记录}
原文来自 MIT OpenCourseWare，作者与课程如上。按来源网站标示的 Creative Commons Attribution--NonCommercial--ShareAlike 4.0 许可复用；本转录及整理沿用同一许可。2026-09-28：助手转录、加入来源定位与简短标注；不暗示作者或 MIT 认可本次整理。许可：\url{https://creativecommons.org/licenses/by-nc-sa/4.0/}。
\end{document}
